\section{Full semantics}
\label{sec:second-order-logic-02}

\begin{definition}[Full second-order assignment]
Let \(\mathcal M\) be a \(\Lang_2\)-structure with domain \(M\).  A full
second-order assignment \(s\) maps each first-order variable to an element of
\(M\) and each \(n\)-ary relation variable to an arbitrary member of
\(\Pow(M^n)\).
\end{definition}

The first-order atomic and connective clauses are unchanged.  The new atomic
and quantifier clauses are
\[
  \mathcal M,s\entails X(t_1,\ldots,t_n)
  \quad\Longleftrightarrow\quad
  \tuple{\sem{t_1}^{\mathcal M}_s,\ldots,
         \sem{t_n}^{\mathcal M}_s}\in s(X),
\]
and
\[
\begin{aligned}
  \mathcal M,s\entails\forall X^n\,\varphi
  \quad\Longleftrightarrow\quad
  &\text{for every }R\subseteq M^n,
    \ \mathcal M,s[X\mapsto R]\entails\varphi,\\
  \mathcal M,s\entails\exists X^n\,\varphi
  \quad\Longleftrightarrow\quad
  &\text{for some }R\subseteq M^n,
    \ \mathcal M,s[X\mapsto R]\entails\varphi.
\end{aligned}
\]

\begin{definition}[Full semantic consequence]
For a set \(\Gamma\) of second-order formulas,
\[
  \Gamma\entails_{\mathrm{full}}\varphi
\]
iff every structure under every full second-order assignment that
satisfies every member of \(\Gamma\) also satisfies \(\varphi\).
\end{definition}

The word \emph{full} modifies the range of second-order quantification.  It
does not mean that the first-order domain contains every set or that the
structure is a set-theoretic universe.

\begin{theorem}[Full comprehension validity]
Let \(\varphi(\bar x)\) be a formula in which \(X^n\) is not free.  Under full
semantics the universal closure of
\[
  \exists X^n\,\forall x_1\cdots\forall x_n\,
  \bigl(X(\bar x)\leftrightarrow\varphi(\bar x)\bigr)
\]
is valid.
\end{theorem}

\begin{proof}
Fix \(\mathcal M\) and a full second-order assignment \(s\) of all
parameters.  Let
\[
  R=\set{\bar a\in M^n:
          \mathcal M,s[\bar x\mapsto\bar a]\entails\varphi(\bar x)}.
\]
Full semantics permits \(X\) to range over this particular subset of \(M^n\).
With \(s(X)=R\), the biconditional follows from the definition of \(R\).
\end{proof}

